\begin{table}[htbp]
\centering
\begin{threeparttable}
\caption{Standardized Effect Sizes}
\label{tab:sde}
\footnotesize
\begin{tabular}{lcccccc}
\toprule
Outcome & $\hat{\beta}$ & SE & SD($Y$) & SDE & SE(SDE) & Classification \\
\midrule
\multicolumn{7}{l}{\textit{Panel A: Pooled}} \\
\addlinespace
Log farm count & -0.0011 & 0.0003 & 1.249 & -0.024 & 0.007 & Small negative \\
Log livestock units & -0.0010 & 0.0003 & 2.064 & -0.013 & 0.004 & Small negative \\
LU per farm & -0.4215 & 0.1014 & 95.692 & -0.124 & 0.030 & Moderate negative \\
\addlinespace
\multicolumn{7}{l}{\textit{Panel B: Heterogeneous (farm count, by livestock composition)}} \\
\addlinespace
Cattle-intensive munis & -0.0035 & 0.0011 & 1.217 & -0.082 & 0.024 & Moderate negative \\
Pig/poultry-intensive munis & -0.0002 & 0.0001 & 0.995 & -0.007 & 0.003 & Small negative \\
\bottomrule
\end{tabular}
\begin{tablenotes}\footnotesize
\item \textit{Notes:} \textbf{Country:} Netherlands. \textbf{Research question:} Does the Dutch piekbelasters (peak emitter) livestock buyout program, the world's largest environmental farm buyout at EUR 1.5 billion, reduce farm counts and livestock in nitrogen-exposed municipalities? \textbf{Policy mechanism:} The Lbv-plus program offers 120\% of assessed farm value to livestock operations near Natura 2000 protected sites that exceed nitrogen emission thresholds, conditional on permanent cessation of agricultural activity at that location. \textbf{Outcome definition:} Log number of agricultural holdings (farm count) and log livestock units (EU-standardized: cattle = 1.0, pig = 0.3, chicken = 0.014) at the municipality-year level from CBS agricultural census. \textbf{Treatment:} Continuous; interaction of Natura 2000 5km buffer share and pre-2019 mean livestock density in thousands of livestock units. \textbf{Data:} CBS 80781ned (agricultural census, 2000--2025), PDOK Natura 2000 boundaries; 328 municipalities, 26 years, 8,053 municipality-year observations. \textbf{Method:} Continuous treatment DiD with municipality and year fixed effects plus municipality-specific linear time trends; standard errors clustered by municipality; robustness via leave-three-out and 500-iteration randomization inference. \textbf{Sample:} All Dutch municipalities with at least one agricultural holding in any year; municipalities with missing Natura 2000 spatial data excluded. SDE $= \hat{\beta} \times \text{SD}(X) / \text{SD}(Y)$ where SD($X$) is the cross-sectional standard deviation of treatment exposure and SD($Y$) is the pre-treatment standard deviation of the outcome. Classification refers to magnitude, not statistical significance: Large ($|$SDE$| > 0.15$), Moderate ($0.05$--$0.15$), Small ($0.005$--$0.05$), Null ($< 0.005$).
\end{tablenotes}
\end{threeparttable}
\end{table}
